\documentclass[10pt,a4paper]{article}

\usepackage[T1]{fontenc}
\usepackage[utf8]{inputenc}
\usepackage[ngerman]{babel}
\usepackage[a4paper,margin=14mm]{geometry}
\usepackage{lmodern}
\usepackage[table]{xcolor}
\usepackage{tabularx}
\usepackage{array}
\usepackage{microtype}
\usepackage[most]{tcolorbox}

\pagestyle{empty}
\setlength{\parindent}{0pt}
\setlength{\parskip}{0pt}
\renewcommand{\familydefault}{\sfdefault}

\definecolor{Navy}{HTML}{17324D}
\definecolor{Blue}{HTML}{3B6EA5}
\definecolor{Sky}{HTML}{EAF3FA}
\definecolor{Mint}{HTML}{E8F5EF}
\definecolor{Gold}{HTML}{F3B33D}
\definecolor{Ink}{HTML}{1E2935}
\definecolor{Muted}{HTML}{617080}
\definecolor{Line}{HTML}{D8E1E8}
\definecolor{CodeBg}{HTML}{F5F7F9}

\color{Ink}

\newcommand{\sectiontitle}[1]{%
  \vspace{2.3mm}
  {\color{Blue}\bfseries\large #1}\par
  \vspace{0.8mm}
  {\color{Blue}\rule{\linewidth}{0.7pt}}\par
  \vspace{1.2mm}
}

\begin{document}

\colorbox{Navy}{%
  \parbox{\dimexpr\linewidth-2\fboxsep\relax}{%
    \vspace{3mm}\color{white}
    \begin{tabularx}{\linewidth}{@{}Xr@{}}
      {\small\bfseries INFORMATIK · KLASSE 10} &
      \colorbox{Gold}{\color{Navy}\small\bfseries\strut LÖSUNGEN}
    \end{tabularx}
    \vspace{2.5mm}

    {\fontsize{22}{25}\selectfont\bfseries Sicherungsblatt: Aufgabe 1}\par
    \vspace{1mm}
    {\large Grundlagen relationaler Datenbanken}
    \vspace{3mm}
  }%
}

\vspace{3mm}
\colorbox{Sky}{%
  \parbox{\dimexpr\linewidth-2\fboxsep\relax}{%
    \vspace{1.8mm}
    \textcolor{Blue}{\bfseries Ziel:}
    Du kennst den Aufbau einer Relation, überträgst Klassen, Attribute und
    Objekte in Tabellenbegriffe, ordnest Datentypen zu und liest ein
    \texttt{CREATE TABLE}-Schema.
    \vspace{1.8mm}
  }%
}

\sectiontitle{1 · Relation = Tabelle}

In einer relationalen Datenbank werden alle Daten in \textbf{Tabellen}
gespeichert. Der Fachbegriff dafür ist \textbf{Relation}.

\vspace{1.5mm}
\textcolor{Blue}{\bfseries Tabelle \texttt{Lehrkraft}:}

\vspace{1mm}
\renewcommand{\arraystretch}{1.2}
\begin{tabularx}{\linewidth}{@{}>{\ttfamily\small}p{16mm}>{\ttfamily\small}p{28mm}>{\ttfamily\small}p{28mm}>{\ttfamily\small}p{16mm}>{\ttfamily\small}X@{}}
  \rowcolor{Sky}\bfseries Kürzel & \bfseries Name & \bfseries Vorname & \bfseries Fach1 & \bfseries Fach2 \\
  bab & Baumbart & Berta & M & Ph \\
  \rowcolor{CodeBg}sch & Schmidt & Peter & D & E \\
\end{tabularx}

\sectiontitle{2 · Von der Klasse zur Tabelle}

\begin{tabularx}{\linewidth}{@{}X X X@{}}
  \textcolor{Blue}{\bfseries Klasse} $\rightarrow$ Tabelle (Relation) &
  \textcolor{Blue}{\bfseries Attribut} $\rightarrow$ Spalte &
  \textcolor{Blue}{\bfseries Objekt} $\rightarrow$ Zeile (Datensatz) \\
\end{tabularx}

\sectiontitle{3 · Datentypen}

\renewcommand{\arraystretch}{1.25}
\begin{tabularx}{\linewidth}{@{}>{\ttfamily\bfseries\small}p{27mm}X@{}}
  INTEGER & ganze Zahlen (z.\,B. 42) \\
  VARCHAR(n) & Text mit variabler Länge, höchstens $n$ Zeichen (z.\,B. \textquotedbl{}Anna-Maria\textquotedbl{}) \\
  CHAR(n) & Text mit fester Länge $n$ (z.\,B. \textquotedbl{}w\textquotedbl{}) \\
  DATE & Datum im Format YYYY-MM-DD (z.\,B. '2008-05-12') \\
\end{tabularx}

\sectiontitle{4 · Primärschlüssel}

Der \textbf{Primärschlüssel} identifiziert jeden Datensatz eindeutig.

\vspace{1.5mm}
\begin{tabularx}{\linewidth}{@{}p{7mm}X@{}}
  \textcolor{Gold}{\Large\bfseries !} & \textbf{UNIQUE:} Jeder Wert kommt in der Spalte nur einmal vor. \\[1mm]
  \textcolor{Gold}{\Large\bfseries !} & \textbf{NOT NULL:} Das Feld darf nie leer sein. \\
\end{tabularx}

\vspace{1.3mm}
Deshalb eignen sich eigens vergebene Nummern wie eine Personalnummer besser
als Namen oder E-Mail-Adressen.

\sectiontitle{5 · CREATE TABLE lesen}

\begin{tabularx}{\linewidth}{@{}p{72mm}X@{}}
  \colorbox{CodeBg}{%
    \parbox{70mm}{%
      \ttfamily\footnotesize
      CREATE TABLE Videospiel (\\
      \ \ Spiel\_ID INTEGER PRIMARY KEY,\\
      \ \ Titel VARCHAR(100),\\
      \ \ Erscheinungsjahr INTEGER,\\
      \ \ Preis\_in\_Cent INTEGER\\
      );
    }%
  }
  &
  \parbox[t]{\linewidth}{%
    Relation: \texttt{Videospiel}\par\vspace{1.5mm}
    Primärschlüssel: \texttt{Spiel\_ID}\par\vspace{1.5mm}
    \texttt{Titel} fasst höchstens 100 Zeichen
  }
  \\
\end{tabularx}

\sectiontitle{6 · Merke}

\begin{tcolorbox}[colback=white,colframe=Line,arc=3mm,boxrule=0.7pt,left=3mm,right=3mm,top=2mm,bottom=2mm]
  - Alle Daten stehen in Tabellen (Relationen) aus Spalten und Zeilen.\par
  - Klasse → Tabelle, Attribut → Spalte, Objekt → Zeile (Datensatz).\par
  - Jede Spalte hat einen Datentyp: \texttt{INTEGER}, \texttt{VARCHAR(n)},
  \texttt{CHAR(n)} oder \texttt{DATE}.\par
  - Der Primärschlüssel ist \texttt{UNIQUE} und \texttt{NOT NULL} und
  identifiziert jeden Datensatz eindeutig.\par
  - \texttt{CREATE TABLE} legt Name, Spalten, Datentypen und Primärschlüssel
  einer Relation fest.
\end{tcolorbox}

\vfill
{\color{Line}\rule{\linewidth}{0.5pt}}\par
\vspace{1.5mm}
\begin{tabularx}{\linewidth}{@{}Xr@{}}
  {\small\color{Muted}Klasse 10 · Datenbanken} &
  {\small\color{Muted}Sicherungsblatt · Aufgabe 1}
\end{tabularx}

\end{document}
